\chapter{일을 나누고 모으기: \texorpdfstring{\code{search}}{search}와 조각 구동}\label{ch:14-rust-search}
\begingroup\setlength{\emergencystretch}{3em}% long identifiers: allow looser lines in this chapter only

\ref{ch:rust-check}장의 마지막 절에서 \code{search}가 일을 나눠 컴퓨터의 처리 장치 코어에 나눠 준다는 것을
보았다. 이 장은 그 둘레를 읽는다. 일이 모두 얼마나 많은지, 어떤 단위로 자르는지, 잘라서 돌린 결과를 어떻게
원래 순서대로 모으는지. 그리고 전시를 위해 더한 두 가지 도구, 느린 드럼 하나를 세워 둔 채 도는 \emph{조각
구동}(\code{fixed\_right})과 정한 선만 골라 메뉴를 만드는 \code{Menu::subset}을 본다. 파일은
\code{search.rs}(126줄)와 \code{bombe.rs}, \code{menu.rs}의 일부다.

\section{일의 양을 센다}\label{sec:s-count}

봄베의 일은 ``드럼 위치 하나에서 전류를 퍼뜨리고 테스트 레지스터를 본다''(\ref{sec:propagate}절)의 반복이다.
몇 번 반복하는지는 서로 무관한 선택들의 곱이다.

\begin{center}
\small
\begin{tabular}{lrrr}
\toprule
기종 & 로터 순서 & 순서당 위치 & 전체 위치 \\
\midrule
Enigma~I & 60 $\times$ 반사판 2 = 120 & $26^3 = 17{,}576$ & 2,109,120 \\
M3 & 336 $\times$ 2 = 672 & $26^3 = 17{,}576$ & 11,811,072 \\
M4 & 336 $\times$ 2 $\times$ 그리스 로터 2 = 1,344 & $26^4 = 456{,}976$ & 614,175,744 \\
\bottomrule
\end{tabular}
\end{center}

로터 순서의 수는 \code{WheelOrder::all}이 만드는 목록의 길이이고(주석에 ``Enigma I: 60 × 2, M3: 336 × 2,
M4: 336 × 2 × 2''라고 적혀 있다), 순서당 위치는 \code{run\_slice}의 네 겹 반복문(코드~\ref{lst:run})이 도는
횟수다. 넘어감 가설(\code{turnover\_hypotheses})을 켜면 여기에 메뉴 구간의 길이가 한 번 더 곱해진다. 가설
``넘어감 없음''에 더해, 메뉴가 쓰는 크립 위치의 가장 작은 값 다음부터 가장 큰 값까지 각 위치마다 가설 하나를
세우기 때문이다.

\begin{lstlisting}[language=Rust,caption={\code{search.rs}: 넘어감 가설의 목록.},label={lst:s-turnovers}]
let mut turnovers: Vec<Option<u16>> = vec![None];
if opts.turnover_hypotheses {
    let lo = menu.edges.iter().map(|e| e.pos).min().unwrap_or(0);
    let hi = menu.edges.iter().map(|e| e.pos).max().unwrap_or(0);
    turnovers.extend((lo + 1..=hi).map(Some));
}
\end{lstlisting}

가설은 \code{None} 하나와 \code{lo+1}부터 \code{hi}까지이므로 모두 $\mathit{hi}-\mathit{lo}+1$개, 곧 메뉴의
구간 길이(\code{Menu::span})다. 넘어감이 \code{lo}에서 일어났다는 가설은 따로 세우지 않는다. 그 경우는 메뉴의
모든 선이 넘어간 뒤의 위치를 쓰므로, 가운데 드럼이 한 칸 다른 ``넘어감 없음''과 같기 때문이다. 노트북의 예에서는
구간이 20인 메뉴로 위치 17,576곳이 351,520곳으로 늘었다.

\Why 위치 하나의 일은 얼마나 큰가. 전류의 확산에서 선 하나는 많아야 한 번 목록에 들어가고, 꺼낼 때마다 그
케이블에 이어진 스크램블러 수만큼 표를 읽는다(\ref{sec:propagate}절). 그러므로 위치 하나의 일은
``선 676가닥 $\times$ (케이블당 연결 수 $+1$)''을 넘지 않는다. 틀린 위치에서는 테스트 레지스터가 금세
가득 차서 그보다 훨씬 일찍 끝난다(조기 종료). 스크램블러 표는 가운데 로터의 반복마다 26개를 새로 만들므로,
3로터 기종의 로터 순서 하나에 표를 만드는 일은 $26^2 = 676$번이고 위치 하나당으로 나누면 표 하나 꼴이다.

\General 다이얼이 세 개 달린 자물쇠는 $10^3$가지를 하나씩 돌려 보면 언젠가 열린다. 다이얼이 하나 늘 때마다
일은 열 배가 된다. 봄베도 같아서, M4의 그리스 로터 하나가 일을 26배로 만든다. 그래도 봄베가 해 볼 만한 것은
가장 큰 다이얼을 돌리지 않기 때문이다. 플러그보드의 경우의 수(케이블 10개일 때 약 $1.5\times10^{14}$,
노트북의 표)는 대입하지 않고 전류의 추론으로 지운다(\ref{sec:bombe-one-page}절).

\begin{figure}[ht]
\centering
\begin{tikzpicture}[x=1cm,y=1cm]
  \draw[cgray,rounded corners=2pt] (0,0) rectangle (12.4,4.6);
  \node[anchor=north west,font=\small\ttfamily] at (0.15,4.5) {for g in greeks};
  \node[anchor=north east,note] at (12.25,4.5) {M4만 26번, 아니면 1번};
  \draw[cgray,rounded corners=2pt] (0.4,0.25) rectangle (12.0,3.95);
  \node[anchor=north west,font=\small\ttfamily] at (0.55,3.85) {for l in 0..26};
  \draw[cblue,rounded corners=2pt,fill=cbluebg!40] (0.8,0.5) rectangle (11.6,3.3);
  \node[anchor=north west,font=\small\ttfamily] at (0.95,3.2) {for m in 0..26};
  \node[blkb,anchor=north west,minimum width=40mm] at (1.0,2.75) {표 26개 만들기\\{\scriptsize \code{Plugged::tables}}};
  \node[note,anchor=north west,text=cblue] at (1.0,1.4) {로터 순서당 676번};
  \draw[corange,rounded corners=2pt,fill=corangebg!50] (5.6,0.7) rectangle (11.4,2.85);
  \node[anchor=north west,font=\small\ttfamily] at (5.75,2.75) {for r in rights};
  \node[blko,anchor=north west,minimum width=50mm] at (5.85,2.3) {\code{propagate} + 판정\\{\scriptsize 로터 순서당 17,576번}};
  \node[note,anchor=north west,text=corange] at (5.85,1.15) {\code{fixed\_right}이면 1번};
\end{tikzpicture}
\caption{\code{run\_slice}의 반복문과 그 안의 일. 표는 가운데 로터가 바뀔 때만 만들고, 확산은 위치마다
한다. 조각 구동은 가장 안쪽 반복을 한 번으로 줄인다(\ref{sec:s-slice}절).}
\label{fig:g-s-loops}
\end{figure}

\section{일을 자르는 단위}\label{sec:s-tasks}

\begin{lstlisting}[language=Rust,caption={\code{search.rs}: 그리스 로터 조각과 일의 목록.},label={lst:s-tasks}]
/// Greek-wheel slices of one wheel order: 26 for M4, a single whole-order slice otherwise.
fn slices(order: &WheelOrder) -> Vec<Option<u8>> {
    if order.greek.is_some() {
        (0..26).map(Some).collect()
    } else {
        vec![None]
    }
}

let tasks: Vec<(usize, Option<u16>, Option<u8>)> = orders
    .iter()
    .enumerate()
    .flat_map(|(i, o)| {
        let turnovers = &turnovers;
        slices(o).into_iter().flat_map(move |g| turnovers.iter().map(move |&t| (i, t, g)))
    })
    .collect();
\end{lstlisting}

일 하나는 (로터 순서의 번호 \code{i}, 넘어감 가설 \code{t}, 그리스 로터 위치 \code{g})의 세 값으로
정해진다. \code{flat\_map}은 ``원소마다 목록을 만들어 이어 붙인다''는 뜻이다. 로터 순서마다 그리스 로터
조각의 목록을 만들고, 조각마다 넘어감 가설의 목록을 만들어 모두 한 줄로 늘어놓는다.

\Why 왜 로터 순서 하나를 통째로 일 하나로 삼지 않는가. 파일 첫머리의 주석이 이유를 적고 있다. ``네 로터 기종의
로터 순서 몇 개만 돌려도 모든 코어가 바쁘도록''. 노트북처럼 M4 로터 순서를 6개만 돌리면 일이 6개뿐이어서,
스레드 24개 가운데 18개가 논다. 그리스 로터 위치로 26조각을 내면 일이 156개가 되어 모든 코어에 고루 돌아간다.
Enigma~I는 로터 순서가 120개여서 자르지 않아도 충분하다. 반대로 너무 잘게 자르면 일을 나눠 주고 거두는 수고가
일 자체보다 커진다. 그리스 로터 한 위치(17,576 위치)는 그 수고를 무시할 만큼 크다.

\General 손님이 여럿인 식탁에서 피자를 자르는 것과 같다. 두 조각으로 자르면 나머지 손님이 기다리고, 백 조각으로
자르면 자르는 데 시간이 다 간다. 일꾼 수보다 넉넉히 많되, 한 조각이 나눠 주는 수고보다 충분히 큰 단위가 좋다.

\Ours \code{run\_slice}의 넷째 인자 \code{greek}가 이 조각을 받는다(\code{Some(g)}이면 그리스 로터 위치
\code{g}의 $26^3$ 위치만 돈다). 같은 함수를 조각 하나에도, 로터 순서 전체에도 쓴다. 넘어감 가설도 조각과 같은
방식으로 일을 늘린다. 가설마다 표가 다르므로 가설끼리도 서로 독립이다.

\section{순서를 지키며 모은다}\label{sec:s-gather}

봄베만 돌리고 체킹은 하지 않는 \code{run\_orders}가 모으는 방법을 잘 보여 준다. Python의
\code{Bombe.run}이 부르는 함수가 이것이다(\ref{ch:15-python-api}장).

\begin{lstlisting}[language=Rust,caption={\code{search.rs}: 로터 순서별 결과로 모은다.},label={lst:s-runorders}]
pub fn run_orders(menu: &Menu, orders: &[WheelOrder], opts: &BombeOptions, parallel: bool) -> Vec<RunResult> {
    let plugged = Plugged::new(menu);
    let tasks: Vec<(usize, Option<u8>)> =
        orders.iter().enumerate().flat_map(|(i, o)| slices(o).into_iter().map(move |g| (i, g))).collect();
    let parts = par_map(&tasks, parallel, |&(i, g)| (i, run_slice(&plugged, &orders[i], opts, g)));
    let mut out: Vec<RunResult> = orders.iter().map(|_| RunResult { stops: Vec::new(), positions: 0 }).collect();
    for (i, r) in parts {
        out[i].positions += r.positions;
        out[i].stops.extend(r.stops);
    }
    out
}
\end{lstlisting}

일은 아무 순서로나 끝난다. 어느 코어가 먼저 끝낼지는 알 수 없다. 그런데 \code{par\_map}의 결과
\code{parts}는 \emph{일의 목록과 같은 순서}로 나온다. 병렬 처리 라이브러리(rayon)의 \code{par\_iter().map(...)
.collect()}는 끝난 순서가 아니라 원래 자리에 결과를 넣기 때문이다. 그래서 4--7행이 조각들을 차례로 이어 붙이면
그리스 로터 0의 정지, 1의 정지, \ldots{} 순으로 쌓이고, 이것은 한 스레드에서 \code{run\_slice(\ldots, None)}을
돌린 결과와 원소 하나하나까지 같다.

\Why 정지의 \emph{집합}만 맞으면 될 것 같지만 이 프로젝트에서는 순서가 중요하다. 기준 RTL과 FPGA 엔진의 시험은
정지 목록을 순서까지 비교한다(\ref{sec:e-verify}절). 소프트웨어가 병렬로 돌았다는 이유로 순서가 흔들리면 그
비교를 할 수 없다. 병렬화는 빠르게 하는 수단일 뿐 답을 바꾸면 안 된다는 원칙을 순서까지 지킨 것이다.

\General 시험지를 여러 채점자가 나눠 매기더라도, 성적표는 번호 순으로 정리해 돌려준다. 누가 먼저 끝냈는지는
성적표에 드러나지 않는다.

\Ours \code{search}는 같은 방식으로 모으되 두 가지를 더한다(코드~\ref{lst:s-search-tail}). 하나는 일 안에서
곧바로 체킹을 하는 것이다(\ref{ch:rust-check}장의 코드). 체킹도 정지마다 독립이므로 같은 일꾼이 이어서 하면 두
번째 분배가 필요 없다. 다른 하나는 해를 독일어 점수 순으로 정렬하는 것이다. 해는 사람이 읽을 것이므로 가장
그럴듯한 것이 먼저 오게 한다.

\begin{lstlisting}[language=Rust,caption={\code{search.rs}: 모은 뒤의 정리.},label={lst:s-search-tail}]
// The bombe cannot tell Enigma I from M3 when only rotors I-V are used; report the model asked for.
if opts.model == Model::M3 {
    for s in &mut report.solutions {
        if s.key.model == Model::EnigmaI {
            s.key.model = Model::M3;
        }
    }
}
report.solutions.sort_by(|a, b| b.score.total_cmp(&a.score));
\end{lstlisting}

M3에 관한 세 줄은 봄베의 원리에서 나온다. 로터 I--V와 반사판 B, C만 쓴 키는 Enigma~I로도 M3로도 똑같이
암호화한다. 체킹은 해를 만들 때 로터의 이름으로 기종을 추정하므로(\code{model\_of}) Enigma~I라고 적는다.
사용자가 M3를 찾았으면 그 이름으로 돌려준다. 계산이 아니라 표기의 문제다.

\section{조각 구동: 느린 드럼을 세워 둔다}\label{sec:s-slice}

\begin{lstlisting}[language=Rust,caption={\code{bombe.rs}: \code{run\_slice}의 가장 안쪽 반복과 그 범위.},label={lst:s-rights}]
let rights = match opts.fixed_right {
    Some(r) => r % N as u8..r % N as u8 + 1,
    None => 0..N as u8,
};
for r in rights {
    positions += 1;
    ...
}
\end{lstlisting}

\code{BombeOptions}에 \code{fixed\_right: Option<u8>}을 두었다. 값이 있으면 오른쪽 로터의 코어 위치를 그 값
하나로 고정하고, 그리스, 왼쪽, 가운데 로터만 돈다. 3로터 기종에서는 $26\times26 = 676$ 위치다. 바뀐 코드는
\code{rights}의 범위를 정하는 네 줄뿐이다.

\paragraph{왜 오른쪽인가.} 모델에서 오른쪽 로터는 가장 안쪽 반복, 곧 가장 빨리 도는 드럼이다(\ref{ch:rust-bombe}장의
``실제 기계와 다른 점''). 그런데 실제 봄베와 이 프로젝트의 두 실물(릴레이식, 기어식)은 에니그마의 \emph{왼쪽}
로터에 해당하는 드럼을 가장 빨리 돌리고, 오른쪽 로터에 해당하는 드럼을 가장 느리게 돌린다. 느린 드럼은
위치 676개마다 한 칸 움직인다. 그러므로 ``오른쪽 코어를 고정한 676 위치''는 실물 기계에서는 \emph{느린 드럼이
한 자리에 서 있는 동안의 한 바퀴}, 전체 구동의 $1/26$이다. 빠른 드럼을 고정할 수는 없다. 빠른 드럼이 한 바퀴
돌아야 가운데 드럼이 한 칸 가는 자리올림의 출발점이기 때문이다(\code{hardware/common/design-brief.md} 5.2절).

\Why 전시에서 이 조각이 필요한 이유는 시간이다. 설계한 릴레이식 봄베는 위치 하나에 63밀리초가 걸려 전체
17,576 위치에 18.5분이 걸리지만, 676 위치는 43초다(같은 문서 5.1절). 관람객이 카드를 고르고 배선하고 정지를
읽는 실습 한 번이 약 7--8분이 되는 것은 이 덕분이다. 조각이 정답 위치를 품도록, 카드가 느린 드럼의 글자를
알려 준다(\ref{ch:32-menu-cards}장).

\begin{figure}[ht]
\centering
\begin{tikzpicture}[x=1cm,y=1cm]
  % pseudo-3D cube: x = middle, y = left, depth = right
  \coordinate (A) at (0,0); \coordinate (B) at (3.0,0); \coordinate (C) at (3.0,3.0); \coordinate (D) at (0,3.0);
  \coordinate (E) at (1.5,0.9); \coordinate (F) at (4.5,0.9); \coordinate (G) at (4.5,3.9); \coordinate (H) at (1.5,3.9);
  \draw[cgray] (A) -- (B) -- (C) -- (D) -- cycle;
  \draw[cgray] (B) -- (F) -- (G) -- (C); \draw[cgray] (D) -- (H) -- (G);
  \draw[cgray,dashed] (A) -- (E) -- (F); \draw[cgray,dashed] (E) -- (H);
  % the slice at depth t
  \fill[corange,opacity=0.35] ($(A)!0.55!(E)$) -- ($(B)!0.55!(F)$) -- ($(C)!0.55!(G)$) -- ($(D)!0.55!(H)$) -- cycle;
  \draw[corange,line width=1pt] ($(A)!0.55!(E)$) -- ($(B)!0.55!(F)$) -- ($(C)!0.55!(G)$) -- ($(D)!0.55!(H)$) -- cycle;
  \node[note,anchor=north] at ($(A)!0.5!(B)$) {가운데 $m$ (26)};
  \node[note,anchor=east] at ($(A)!0.5!(D)$) {왼쪽 $l$ (26)};
  \node[note,anchor=north west] at ($(B)!0.5!(F)$) {오른쪽 $r$ (26)};
  \node[note,text=corange,anchor=west] at ($(C)!0.55!(G)+(0.2,0.1)$) {$r$ 고정: $26\times26=676$};
  \begin{scope}[xshift=8.2cm]
    \node[blkn,minimum width=44mm,align=left,anchor=west] (full) at (0,2.4) {전체 구동: 17,576 위치\\{\scriptsize 릴레이식 63\,ms/위치이면 18.5분}};
    \node[blko,minimum width=44mm,align=left,anchor=west] (sl) at (0,0.9) {조각 구동: 676 위치\\{\scriptsize 같은 속도로 43초}};
  \end{scope}
\end{tikzpicture}
\caption{드럼 위치의 공간 $26^3$과 조각 하나. 실물 기계에서 오른쪽 로터의 드럼은 가장 느린 드럼이므로, 조각은
그 드럼이 한 자리에 서 있는 동안의 구동이다. 시간은 \code{hardware/common/design-brief.md} 5.1절의 설계값이다.}
\label{fig:g-s-slice}
\end{figure}

\paragraph{조각의 답은 전체의 답의 일부다.} 조각 구동의 정지 목록은 전체 구동의 정지 목록에서 오른쪽 코어가
그 값인 것만 골라낸 것과 \emph{순서까지} 같다. 시험이 정확히 이것을 확인한다.

\begin{lstlisting}[language=Rust,caption={\code{tests/search.rs}: 조각 구동의 시험(일부).},label={lst:s-slice-test}]
let full = run(&sub, &order, &BombeOptions::default());
let slice = run(&sub, &order, &BombeOptions { fixed_right: Some(c[2]), ..Default::default() });
assert_eq!(slice.positions, 676);
let expected: Vec<_> = full.stops.iter().filter(|s| s.position.right == c[2]).cloned().collect();
assert_eq!(slice.stops, expected);
assert!(slice.stops.iter().any(|s| [s.position.left, s.position.middle] == [c[0], c[1]]));
\end{lstlisting}

\Why 이것이 성립하는 근거는 위치끼리 아무것도 주고받지 않는다는 데 있다. \code{propagate}는 위치마다 상태를
0에서 새로 시작하고(\code{let mut live = [0u32; N]}), 그 위치의 표만 읽는다. 앞 위치의 결과가 다음 위치에 남지
않으므로 위치 몇 개를 건너뛰어도 나머지 위치의 판정은 바뀌지 않는다. 반복의 순서는 그대로이므로 남은 정지의
순서도 그대로다.

\General 반 전체의 시험지 가운데 3반 학생 것만 채점해도 3반 학생의 점수는 전체를 채점할 때와 같다. 시험지끼리
서로 영향을 주지 않기 때문이다. 반대로 상대 평가였다면 일부만 채점한 점수는 달라진다.

\Ours 조각 구동에서는 \ref{sec:s-count}절의 ``표 하나를 26 위치가 나눠 쓴다''가 사라진다. 가운데 로터마다
표 26개를 만들지만 그 표를 쓰는 위치는 하나뿐이다(선마다 $(r+k) \bmod 26$번째 표를 읽으므로, 읽는 표는 많아야
선의 수만큼이다). 위치당 일은 늘지만 위치가 676개뿐이어서 상관없다. 빠르게 하는 것보다
코드를 하나로 두는 쪽을 택했다. FPGA 엔진의 레지스터 맵에는 이 선택이 없다(표~\ref{tab:g-e-regmap}). 엔진은
17,576 위치를 몇 밀리초에 돌므로 조각으로 나눌 이유가 없고, 조각 구동은 실물 두 대를 위한 것이다.

\section{정한 선으로 메뉴 만들기: \texorpdfstring{\code{Menu::subset}}{Menu::subset}}\label{sec:s-subset}

\ref{ch:rust-menu}장의 \code{select}는 선을 탐욕법으로 \emph{골라 준다}. 전시 카드처럼 어떤 선을 쓸지 이미
정해져 있을 때는 그 선들을 \emph{지정해서} 메뉴를 만들어야 한다. 그 일을 하는 것이 \code{subset}이다.

\begin{lstlisting}[language=Rust,caption={\code{menu.rs}: 크립 위치로 고른 부분 메뉴.},label={lst:s-subset}]
pub fn subset(&self, positions: &[u16]) -> Result<Menu> {
    let edges: Vec<Edge> = self.edges.iter().copied().filter(|e| positions.contains(&e.pos)).collect();
    if let Some(&p) = positions.iter().find(|&&p| !self.edges.iter().any(|e| e.pos == p)) {
        return Err(Error::Config(format!("crib position {p} is not in the menu")));
    }
    if edges.is_empty() {
        return Err(Error::Config("empty sub-menu".into()));
    }
    let mut menu = Menu { offset: self.offset, crib: self.crib.clone(), edges, test_letter: 0, input_wire: 0 };
    menu.choose_test_letter();
    Ok(menu)
}
\end{lstlisting}

선은 크립 위치(\code{pos})로 지정한다. 한 크립 위치에는 선이 하나뿐이므로 위치가 선의 이름 노릇을 한다.
2행은 지정한 위치의 선만 남기고, 3--5행은 메뉴에 없는 위치를 지정하면 오류로 돌려준다. 조용히 무시하면 사용자가
선이 하나 빠진 메뉴를 그런 줄 모르고 돌리게 된다. 9행이 이 함수의 요점이다.

\Why 부분 메뉴가 원래의 크립(\code{crib})과 그 위치(\code{offset})를 그대로 가져가는 이유는 체킹 때문이다.
체킹의 2단계는 메뉴의 선만이 아니라 \emph{크립의 모든 글자}를 암호문과 대조한다(\ref{ch:rust-check}장).
선 6개짜리 부분 메뉴로 봄베를 돌렸더라도 정지를 검증할 때는 나머지 크립 글자가 모두 쓸모 있다. 반대로 손으로
선을 적어 만드는 \code{Menu::from\_edges}는 크립을 모르므로 \code{offset}을 0, 크립을 빈 글자로 채워 두고,
그런 메뉴의 정지는 체킹할 수 없다.

\Ours 테스트 글자는 부분 메뉴에서 다시 고른다(\code{choose\_test\_letter}: 가장 큰 덩어리 안에서 연결이 가장
많은 글자, 입력 선은 \texttt{A}이되 테스트 글자가 \texttt{A}이면 \texttt{B}). 원래 메뉴의 테스트 글자가 부분
메뉴에서는 연결이 하나뿐이거나 아예 없을 수 있기 때문이다. 테스트 글자를 사람이 따로 정하는 일은 Python 쪽의
\code{Menu.with\_test}가 맡는다(\ref{ch:15-python-api}장). \code{subset}을 쓰는 곳은 둘이다. 조각 구동의
시험(코드~\ref{lst:s-slice-test})이 선 10개 메뉴에서 앞의 8개를 골라 쓰고, 메뉴 카드 생성기가 미리 정한 크립
위치들로 카드의 메뉴를 만든다(\ref{ch:32-menu-cards}장).

\section{정지 하나는 어떻게 검증되는가}\label{sec:s-check}

조각 구동과 부분 메뉴는 \emph{어느 위치를 돌지}와 \emph{어느 선을 걸지}만 바꾼다. 위치 하나를 판정하고
검증하는 방법은 그대로다. 체킹의 1단계는 정지의 위치 하나에서 전류를 다시 흘리는데, 이때 넘기는 선택이
다음과 같다.

\begin{lstlisting}[language=Rust,numbers=none]
let opts = BombeOptions { diagonal_board: diagonal, turnover_at: stop.turnover_at, fixed_right: None };
\end{lstlisting}

\code{fixed\_right}는 \code{None}이다. 체킹은 반복문을 돌지 않고 \code{trace}로 위치 하나만 보므로 조각 구동인지
아닌지가 의미가 없다. 반면 넘어감 가설은 정지에 실려 온 값을 그대로 쓴다. 가설에 따라 선마다 쓰는 표가 달라지기
때문이다. 정지(\code{Stop})가 위치와 테스트 레지스터 말고도 로터 순서와 넘어감 가설을 함께 들고 다니는 이유가
이것이다. 정지 하나만 있으면 그 위치를 다시 만들 수 있다.

\begin{keybox}[title={이 장의 정리}]
\begin{itemize}
\item 일의 양은 서로 독립인 선택의 곱이다. 로터 순서 $\times$ 위치 $\times$ 넘어감 가설. 플러그보드는 곱하지 않는다.
\item 일은 코어 수보다 넉넉히 많게, 그러나 나눠 주는 수고보다 크게 자른다. M4는 그리스 로터 위치로 26조각을 낸다.
\item 병렬로 돌려도 결과는 원래 순서로 모은다. 다른 구현과 정지 목록을 순서까지 비교하기 위해서다.
\item 위치끼리 상태를 주고받지 않으므로, 느린 드럼 하나를 세운 조각 구동의 정지는 전체 구동의 정지를 거른 것과 같다.
\item 부분 메뉴는 크립을 그대로 가져가서 체킹할 수 있다.
\end{itemize}
\end{keybox}
\endgroup
